\subsection{EXPANSIONS TO BASE A}

The most important base sequences are those for which \(d_{n}=a^{n}\) , where a is an integer \textgreater1; a is then said to be the base number (or simply the base) of the corresponding expansions. For numerical calculations, expansions to base 10, which are called decimal expansions, are used; also expansions to base 2 (dyadic expansions) and base 3 (triadic expansions) are often used.

To represent the approximations by defect \(r_n\) to a number \(x \geqslant 0\) , in its expansion to base \(a\) , the following symbolism is employed: each integer \(u\) such that \(0 \leqslant u \leqslant a - 1\) is denoted by a particular sign; if

\[
r _ {n} = p _ {0} + \sum_ {k = 1} ^ {n} \frac {u _ {k}}{d _ {k}},
\]

we first write down, with the aid of these signs, the representation to base \(a\) of the integer \(p_{0}=[x]\geqslant0\) (Set Theory, Chapter III, § 5, no. 7), then we put a point (``decimal point'') and we write

successively the signs representing the numbers \(u_1, u_2, \ldots, u_n\) . If \(S\) is the symbol thus obtained, it is customary to write \(x = S \ldots\) , by abuse of language. It should be understood once and for all that such a relation is only an abbreviated method of indicating that the right-hand side is the approximation to \(x\) to within \(1/a^n\) by defect.

For negative numbers the established usage is different: we write an approximation to \(x' = -x > 0\) in the symbolism described above, and precede it by the sign ``-''; it is thus an approximation to \(x\) to within \(1/a^n\) by excess that is so denoted.

This usage has its inconveniences for numerical calculations. In the notation for negative logarithms the same symbolism is adopted as for positive numbers, by putting a bar over the integral part, to indicate that it is equal to the negative of the number written.

